\section{Discussion}
\label{sec:discussion}

\subsection{Implications}

The broader lesson is that narrow programmable hooks help when applications
vary policy but the kernel should retain mechanism ownership. \namei applies
that lesson to VFS name resolution rather than to a new filesystem. For
\namei, the retained mechanism is VFS path walking and lower-filesystem object
semantics. The \namei boundary is useful only when the correctness condition is
pathname-to-object selection or directory visibility over existing objects.

FUSE, stackable filesystems, and custom filesystems remain the appropriate implementation boundary
when a workload needs filesystem-service features such as synthetic contents,
persistent custom metadata, data-path mediation, write-conflict resolution,
distributed metadata, or cross-path transactions. \namei is intended to avoid
that boundary only for the subset of policies whose effect is visible during
lookup or directory enumeration.

\subsection{Extension Path}

The same boundary also defines future extensions. Synthetic aliases require
additional kernel support, while synthetic contents, custom metadata
persistence, distributed indexing, and cross-path transactions remain filesystem
or runtime responsibilities. \namei is intended only for workloads whose
correctness oracle is path selection or visibility over existing lower objects.
